% Переработанное заключение без перечисления экспериментальных показателей.
% Заголовок \chapter*{Заключение} при необходимости добавляется в основном файле.

В работе решена научно-прикладная задача повышения эффективности оптимизации SQL-запросов за счёт разработки методов, дополняющих штатный стоимостный оптимизатор без его полной замены. Предложены три взаимодополняющих направления: адаптивная настройка параметров стоимостной модели по телеметрии рабочей нагрузки, рекомендация подсказок оптимизатору на основе сопоставления планов выполнения в пространстве векторных представлений и агентный итеративный поиск подсказок при ограниченном бюджете пробных запусков. Поставленная цель достигнута, сформулированные задачи выполнены.

Предложенные методы предназначены для различных практических ситуаций и не должны рассматриваться как обязательные последовательные этапы единого алгоритма. Метод адаптивной настройки воздействует на параметры стоимостной модели, метод рекомендации использует ранее накопленный опыт для быстрого выбора подсказок в онлайн-режиме, а агентный метод решает отдельную задачу оффлайн-поиска новой конфигурации. Совместное применение этих методов в работе экспериментально не исследовалось, поэтому при их комбинировании необходимо отдельно оценивать итоговый план: влияние изменения параметров стоимостной модели и подсказок оптимизатору может быть неаддитивным.

Практическое применение метода адаптивной настройки стоимостной модели целесообразно в случаях, когда оценки кардинальности в целом достаточно корректны, однако штатные параметры стоимости не отражают фактические характеристики аппаратной среды, состояние буферного кэша или особенности рабочей нагрузки. Метод не предназначен для исправления ошибок оценки кардинальности. Поэтому его применению должна предшествовать диагностика причины неэффективного выбора плана. Если основная ошибка возникает на уровне кардинальностей, корректировка стоимостных коэффициентов не устраняет первопричину; если же рассогласование связано со стоимостью ввода-вывода или вычислительных операций, адаптивная настройка может использоваться для уточнения штатной модели.

Состояние буферного кэша в предложенном методе оценивается агрегированно на уровне отношений, без отслеживания отдельных страниц. Это снижает требования к памяти и вычислительные накладные расходы, но требует накопления достаточного числа наблюдений, поэтому результаты начального периода работы следует интерпретировать осторожно. При практическом внедрении необходимо также учитывать взаимодействие динамической оценки резидентности со штатными механизмами СУБД, в частности с параметрами, уже учитывающими предполагаемое кэширование данных. Отдельной задачей дальнейшего развития является исключение возможного частичного двойного учёта эффекта буферного кэша и более тесная интеграция динамической оценки резидентности в штатную стоимостную модель.

Метод рекомендации подсказок на основе сопоставления планов наиболее естественно применять для стационарных или частично повторяющихся аналитических нагрузок, где уже накоплена история выполнения запросов и для нового запроса можно найти структурно близкие ранее наблюдавшиеся планы. Близость векторных представлений используется только для формирования конфигурации-кандидата и не является достаточным основанием для её безусловного применения. После повторного планирования выполняется дополнительная проверка согласованности по накопленному репозиторию. Если подходящего локального опыта недостаточно либо кандидат не проходит проверку, онлайн-метод должен отказаться от рекомендации и сохранить штатную конфигурацию оптимизатора. Такой отказ является нормальным результатом работы метода, а оффлайн-поиск не запускается автоматически как его внутренняя ветвь.

При использовании онлайн-метода следует учитывать стоимость дополнительного планирования, построения векторных представлений и поиска ближайших соседей. Поэтому он прежде всего ориентирован на аналитические запросы, для которых ожидаемый выигрыш от улучшенного плана превышает стоимость этапа рекомендации. При росте репозитория накопленного опыта целесообразно использовать индексированные методы приближённого поиска ближайших соседей вместо полного линейного перебора. Перенос метода на СУБД с существенно иной системой физических операторов или иным представлением планов требует отдельной проверки пригодности используемых векторных представлений.

Агентный метод итеративного поиска подсказок предназначен для оффлайн-оптимизации запросов, когда допустим заранее заданный бюджет пробных запусков. Он может применяться при отсутствии достаточного накопленного опыта, при невозможности сформировать надёжную онлайн-рекомендацию или как самостоятельный механизм дополнительной оптимизации длительно выполняющихся запросов. Языковая модель в этом случае используется как компонент принятия решений, который анализирует текущее состояние поиска, результаты предыдущих попыток и выбирает следующее действие. Метод следует рассматривать как эвристический механизм направленного поиска, ориентированный на рациональное использование ограниченного бюджета, а не как алгоритм, гарантирующий нахождение глобально оптимальной конфигурации.

Практическая ценность агентного подхода возрастает по мере увеличения пространства допустимых конфигураций, когда полный перебор становится слишком дорогим. Вместе с тем при реальном внедрении необходимо учитывать полную стоимость оффлайн-оптимизации, включая не только пробные выполнения запросов, но и время обращения к языковой модели. Применение агентного поиска оправдано прежде всего для длительных или многократно повторяющихся запросов, для которых затраты предварительной оптимизации могут быть компенсированы последующими выполнениями.

Для практического выбора метода администратором базы данных можно использовать следующую схему:
\begin{enumerate}
    \item Если проблема связана преимущественно с несоответствием параметров стоимостной модели фактической среде выполнения при достаточно корректных оценках кардинальности, целесообразно применять адаптивную настройку стоимостной модели.
    \item Если для текущего запроса имеется накопленная история структурно близких планов с известными эффективными конфигурациями подсказок, целесообразно использовать онлайн-метод рекомендации на основе сопоставления планов.
    \item Если близкие аналоги отсутствуют либо найденная рекомендация не проходит проверку, для длительного запроса может быть отдельно запущен агентный оффлайн-поиск с заранее установленным бюджетом пробных запусков.
    \item Результаты оффлайн-оптимизации целесообразно сохранять в репозитории накопленного опыта и использовать при последующих онлайн-рекомендациях, постепенно расширяя область запросов, для которых дорогостоящий поиск не требуется.
    \item При совместном применении нескольких предложенных методов итоговый план необходимо валидировать отдельно, поскольку экспериментальных оснований считать их эффекты независимыми или аддитивными в работе не получено.
\end{enumerate}

Для инженеров СУБД полученные результаты представляют интерес как основа для совершенствования внутренних механизмов планирования. Сопоставление исходных и улучшенных планов позволяет выявлять случаи, в которых штатная стоимостная модель, ограничения пространства поиска или выбор физических операторов приводят к менее эффективному решению. Предложенные методы сохраняют существующий оптимизатор в качестве основного механизма построения планов и дополняют его средствами адаптации, переноса накопленного опыта и направленного поиска.


\medskip
\noindent
\textbf{Перспективы дальнейшего развития направления.}
Дальнейшее развитие работы связано как с углублением отдельных предложенных методов, так и с исследованием их совместного применения в более общей системе поддержки решений по оптимизации запросов.

Для метода адаптивной настройки стоимостной модели перспективным направлением является более тесная интеграция динамической оценки резидентности отношений со штатными механизмами СУБД. В частности, представляет интерес использование наблюдаемого состояния буферного кэша непосредственно при оценке ожидаемого числа обращений к памяти и диску, что позволит уменьшить риск частичного двойного учёта эффекта кэширования. Отдельного исследования требуют устойчивость метода при длительной конкурентной нагрузке, сценарии с операциями записи, полные накладные расходы на сбор телеметрии и минимальный объём наблюдений, необходимый для выхода модели в устойчивый режим. Перспективной является также разработка диагностического механизма, позволяющего различать ошибки стоимостной модели и ошибки оценки кардинальности и тем самым выбирать корректный способ воздействия на оптимизатор.

Для метода рекомендации подсказок на основе сопоставления планов основным направлением является повышение качества переноса на ранее не наблюдавшиеся шаблоны запросов. Для этого могут исследоваться доменная адаптация моделей векторных представлений, дополнительное обучение на планах выполнения, объединение текстовых и структурных признаков и гибридные метрики близости. Практически важны также разработка явного критерия уверенности рекомендации, масштабирование репозитория с использованием индексных методов приближённого поиска ближайших соседей и учёт устаревания накопленного опыта при изменении данных, конфигурации СУБД и аппаратной платформы. Отдельного изучения требует перенос метода на СУБД с иной системой физических операторов и другим представлением планов.

Для агентного метода наиболее существенным продолжением является переход к более крупным пространствам действий, в которых подсказки задаются не только глобально для классов операторов, но и локально для отдельных соединений, отношений, индексов и способов доступа. В такой постановке полный перебор становится практически невозможным, а направленный поиск при ограниченном бюджете приобретает ещё большее значение. Дополнительным направлением является развитие политики агента за счёт использования накопленных в эксплуатации успешных и неуспешных траекторий для периодического дообучения и формирования более информативного контекста.

Отдельным направлением дальнейшей работы является экспериментальная проверка совместного применения методов глав~2--4. Наиболее общей постановкой здесь является создание диагностического механизма, который по характеристикам запроса, состоянию СУБД и накопленной истории определяет причину неэффективного планирования и выбирает подходящее действие: корректировку стоимостной модели, онлайн-перенос конфигурации по аналогии, агентный оффлайн-поиск либо сохранение штатного решения. Такая система должна учитывать возможное неаддитивное взаимодействие методов и отдельно проверять качество итогового плана после каждого изменения.

Таким образом, диссертационная работа показывает возможность повышения качества планирования аналитических SQL-запросов без полной замены существующего оптимизатора. Предложенные методы могут использоваться как самостоятельные средства в различных диагностических ситуациях и как элементы более общей системы управления качеством планирования. Выбор конкретного метода должен определяться причиной неэффективного планирования, характером рабочей нагрузки, наличием накопленного опыта и допустимой стоимостью дополнительной оптимизации.
